\section{Introduction}
\label{sec:introduction}

For a subset $S$ of a von Neumann algebra $M$, write $W^*(S)$ for the
smallest unital von Neumann subalgebra containing $S$. Thus adjoints
are included and generation is taken in the weak operator topology.
The generator problem asks whether every von Neumann algebra with
separable predual has the form $W^*(x)$ for one bounded operator $x$.
Equivalently, it asks for two self-adjoint generators: their sum
$a+ib$ recovers $a$ and $b$ as its real and imaginary parts.

We prove this assertion for type $\mathrm{II}_1$ factors by first
establishing a relative generation theorem. For a finite factor with
normalized trace $\tau$, put $\norm{x}_2=\tau(x^*x)^{1/2}$.
An inclusion $P\subset M$ is \emph{irreducible} when
$P'\cap M=\mathbb C1$. Write $\mathcal U(M)$ for the unitary group
of $M$.

\begin{theorem}\label{thm:relative-generation}
Let $P\subset M$ be an irreducible inclusion of type $\mathrm{II}_1$
factors, with $M$ of separable predual. The set
\[
 \{u\in\mathcal U(M):W^*(P,u)=M\}
\]
is a dense $G_\delta$ subset of the unitary group $\mathcal U(M)$
for the trace $2$-norm topology.
\end{theorem}

\begin{theorem}\label{thm:main}
Every type $\mathrm{II}_1$ factor $M$ with separable predual is
generated by two self-adjoint elements. Equivalently, there is an
$x\in M$ such that $M=W^*(x)$.
\end{theorem}

Throughout the paper,
\emph{separable} for a finite von Neumann algebra means that its
predual is separable. This ensures separability of $L^2(M,\tau)$;
it does not require operator-norm separability of $M$.

\paragraph{Earlier work.}
The type~I, hyperfinite, and properly infinite cases of the generator
problem were established by Pearcy, Suzuki and Sait\^o, and Wogen,
respectively~\cite{Pearcy1962,SuzukiSaito1963,Wogen1969};
see also the historical discussion in~\cite[Section~3]{Sherman2012}.
Willig's direct-integral theorem reduces the general question to
type~$\mathrm{II}_1$ factors~\cite[Corollary~2]{Willig1974}.
Consequently, Theorem~\ref{thm:main} also gives single generation of
every von Neumann algebra with separable predual. This last consequence
uses the established direct-integral reduction, whereas our proof
below concerns finite factors.

Much of the progress in the finite case has used additional structure.
Behncke established single generation for tensor products of two
$\mathrm{II}_1$ factors~\cite{Behncke1972}.
Popa treated factors with Cartan subalgebras~\cite{Popa1985Cartan};
Ge and Popa treated factors with property~$\Gamma$ and recovered the
tensor-product case~\cite{GePopa1998}; and Ge and Shen obtained
arithmetic group examples with a small-support
refinement~\cite{GeShen2002}.
Shen's generator invariant unified these and other positive
classes~\cite{Shen2009}. Dykema, Sinclair, Smith, and White established
its exact amplification formula and sharpened its connection with the
number of self-adjoint generators~\cite{DSSW2008}.

Sherman showed, in particular, that finite generation of every
countably generated $\mathrm{II}_1$ factor would suffice for a positive
answer~\cite[Theorem~3.8(3)]{Sherman2012}.
More recently, Gao, Kunnawalkam Elayavalli, Patchell, and Tan proved
single generation when a countable unitary generating family lies in
one sequential-commutation orbit in an ultrapower: successive
members can be joined there by finite chains with Haar intermediate
unitaries and commuting adjacent terms~\cite[Theorem~A]{GaoEtAl2025}. Popa's work on tight
decompositions studies stable single generation, which requires every
amplification of a factor to be singly generated~\cite{Popa2021Tight}.

Requiring generators to vary continuously with a family of factors
leads to a different problem. Hayes constructed a $W^*$-bundle---a
continuous family of tracial von Neumann algebras---over a compact
metrizable base, with separable-predual $\mathrm{II}_1$ fibers, for
which no finite family of continuous sections generates every
fiber~\cite[Theorem~1.1]{Hayes2026Bundles}.

Relative generation by one unitary is already available for finite-index
inclusions, by the downward basic construction used in
\cite[Lemma~6.2]{DSSW2008}: the associated Jones projection $e$
generates the extension over the subfactor, and $1-2e$ is unitary.
Theorem~\ref{thm:relative-generation} applies to every irreducible
inclusion and also establishes genericity of the relative generators.

The universal conclusion of Theorem~\ref{thm:main} also makes the
generator invariants $G$ and $G_{\mathrm{sa}}$ vanish for every
separable-predual $\mathrm{II}_1$ factor. It implies that Voiculescu's
microstates free entropy dimensions $\delta$ and $\delta_0$ can change
under replacement of a finite self-adjoint von Neumann generating
tuple. Section~\ref{sec:consequences} gives these deductions and their
precise scope.

\paragraph{The main estimate and the proof.}
Our argument combines Popa's free-independence construction in tracial
ultraproducts~\cite{Popa1995,Popa2014} with a perturbation whose effect
accumulates along a long word. Given a unitary $u$ and a target unitary
$z$, we coordinate a perturbation of $u$ across its repeated occurrences
in a word with coefficients from $P$. With a small parameter $\alpha$
fixed, the perturbation has $2$-norm of order $\alpha/\sqrt\ell$,
where $\ell$ is the number of occurrences. The resulting word has
correlation with a test vector $\xi$ equal to
$(\alpha/2)\tau(\xi^*z)$ up to an error of order
$\alpha^2\norm{\xi}_2+1/\ell$.

The essential estimate is uniform in $\ell$. An operator-norm bound
on each perturbation separately would grow with the length of the
product. We instead dominate all coefficients of the nonlinear terms
by a free-group convolution operator with nonnegative coefficients. The length-one case of
Haagerup's inequality~\cite[Lemma~1.4]{Haagerup1979} bounds its norm
independently of $\ell$. This controls the full exponential remainder,
including cancellations between words. Proposition~\ref{prop:perturbation}
states the precise bounds. The witness word need only retain a fixed
positive correlation as $\ell\to\infty$.
The free variables and perturbation are first constructed in a tracial
ultrapower. For each fixed word length, choosing one coordinate of
their representing sequences gives the required unitaries in $M$ and
$P$, with an arbitrarily small additional error.

For a fixed vector $\xi\in L^2(M,\tau)$, the norm of its projection
onto $L^2(W^*(P,u),\tau)$ is a bounded lower semicontinuous function
of $u$. Baire's theorem supplies common continuity points for a
countable dense family of vectors, and contractivity extends continuity
to every vector. If $W^*(P,u)$ were proper at such a point, a nonzero
bounded vector orthogonal to it would have projection norm zero.
The perturbation estimate gives nearby unitaries for which that
projection norm has a fixed positive lower bound, a contradiction.

Finally, Popa's irreducible hyperfinite embedding theorem
\cite{Popa1981,Popa2021Ergodic} lets us choose $P$ generated by two
self-adjoint elements. Relative generation then gives four
self-adjoint generators for every separable-predual $\mathrm{II}_1$
factor. Applying this uniform bound to a half-trace corner, and using
an explicit two-by-two matrix encoding, gives
Theorem~\ref{thm:main}. This encoding belongs to the established
matrix-generation techniques of Douglas--Pearcy and Wogen
\cite{DouglasPearcy1969,Wogen1969}; we use the von Neumann
algebra form of the matrix encoding
in~\cite[Proposition~1.1]{DSSW2008}.

Section~\ref{sec:independence} states the exact independence input and
constructs the free Haar unitaries required in the ultrapower.
Section~\ref{sec:perturbation} proves the uniform perturbation estimate.
Section~\ref{sec:relative-generation} carries out the continuity
argument and identifies the full relative-generator locus as
$G_\delta$. Section~\ref{sec:two-generators} gives the corner and
matrix construction proving single generation.
Section~\ref{sec:consequences} treats the general von Neumann
algebra, generator-invariant, and free-entropy consequences.

\paragraph{Notation.}
We write $\norm{x}$ for the operator norm and use the inner product
$\langle a,b\rangle=\tau(b^*a)$ on $L^2(M,\tau)$, linear in the
first variable. The same trace notation is used on tracial ultrapowers.
All subalgebras are unital unless a corner unit is specified.
For a von Neumann subalgebra $N\subset M$, $E_N$ denotes the
trace-preserving conditional expectation and $e_N$ the orthogonal
projection from $L^2(M,\tau)$ onto $L^2(N,\tau)$.
